% !TeX program = lualatex
\documentclass[11pt,a4paper]{article}

\usepackage[a4paper,left=18mm,right=18mm,top=17mm,bottom=18mm,footskip=9mm]{geometry}
\usepackage{fontspec}
\usepackage[UTF8,scheme=plain,fontset=none]{ctex}
\setCJKmainfont{FandolHei-Regular.otf}[BoldFont=FandolHei-Bold.otf]
\setmainfont{texgyreheros}[Extension=.otf,UprightFont=*-regular,BoldFont=*-bold,ItalicFont=*-italic,BoldItalicFont=*-bolditalic]
\setsansfont{texgyreheros}[Extension=.otf,UprightFont=*-regular,BoldFont=*-bold,ItalicFont=*-italic,BoldItalicFont=*-bolditalic]
\usepackage{microtype}
\usepackage{xcolor}
\usepackage{enumitem}
\usepackage[explicit]{titlesec}
\usepackage{fancyhdr}
\usepackage{tabularx}
\usepackage{array}
\usepackage{ragged2e}
\usepackage[hidelinks]{hyperref}

\definecolor{accent}{HTML}{315F75}
\definecolor{ink}{HTML}{303438}
\definecolor{muted}{HTML}{666B70}
\definecolor{rulegray}{HTML}{BEC4C8}

\hypersetup{
  colorlinks=true,
  urlcolor=accent,
  linkcolor=accent
}

\pagestyle{fancy}
\fancyhf{}
\renewcommand{\headrulewidth}{0pt}
\fancyfoot[R]{\fontsize{8}{10}\selectfont\color{muted}\thepage}
\renewcommand{\normalsize}{\fontsize{10.5}{14.5}\selectfont}
\AtBeginDocument{\normalsize\color{ink}}
\raggedbottom
\setlength{\parindent}{0pt}
\setlength{\parskip}{0pt}
\linespread{1}
\emergencystretch=1.5em
\setlist[itemize]{leftmargin=1.3em,itemsep=0.24em,topsep=0.22em,parsep=0pt,partopsep=0pt}
\setlist[enumerate]{leftmargin=2.05em,itemsep=0.45em,topsep=0.28em,parsep=0pt,partopsep=0pt}
\renewcommand{\arraystretch}{1.08}
\RaggedRight

\titleformat{\section}
  {\fontsize{13}{16}\selectfont\bfseries\scshape\color{ink}}
  {}
  {0pt}
  {#1\nobreak\hspace{0.65em}{\color{rulegray}\leaders\hrule height 0.45pt\hfill\kern0pt}}
\titlespacing*{\section}{0pt}{18pt}{9pt}

\titleformat{\subsection}
  {\fontsize{12}{15}\selectfont\bfseries\scshape\color{accent}}
  {}
  {0pt}
  {#1}
\titlespacing*{\subsection}{0pt}{10pt}{5pt}

\newcommand{\rightmeta}[1]{{\color{muted}\fontsize{9}{12}\selectfont\mbox{#1}}}

\newcommand{\alignedline}[2]{%
  \noindent\begin{tabular*}{\textwidth}{@{}l@{\extracolsep{\fill}}r@{}}
    #1 & \rightmeta{#2}
  \end{tabular*}%
}

\newcommand{\datedline}[2]{%
  \if\relax\detokenize{#2}\relax
    #1\par
  \else
    \alignedline{#1}{#2}\par
  \fi
}

\newcommand{\roleline}[2]{%
  \alignedline{\itshape #1}{#2}\par
}

\newcommand{\institutionheader}[1]{%
  \par\noindent{\fontsize{9.5}{12}\selectfont\bfseries\color{muted}#1}\par\vspace{5pt}
}

\newcommand{\corrmark}{\textsuperscript{\textcolor{accent}{*}}}
\newcommand{\corrauth}[1]{#1\corrmark}

\newcommand{\educationentry}[5]{%
  \begin{tabularx}{\textwidth}{@{}>{\raggedright\arraybackslash}X r@{}}
    \textbf{#1} & \rightmeta{#5}\\[2pt]
    #2 & \rightmeta{#4}\\[2pt]
    \multicolumn{2}{@{}l@{}}{{\fontsize{9.5}{13}\selectfont\color{muted}#3}}
  \end{tabularx}
}

\newcommand{\pubtopic}[2]{%
  \subsection*{#1}
  {\fontsize{10}{14}\selectfont #2\par}\vspace{8pt}
}

\newcommand{\pub}[1]{\item #1}

\newsavebox{\publinks}
\newcommand{\pubentry}[6][]{%
  \par\noindent
  \begin{minipage}{\textwidth}
    \RaggedRight\hyphenpenalty=10000
    {\fontsize{10}{13}\selectfont\textbf{#2}\par}\vspace{3pt}
    \if\relax\detokenize{#3}\relax\else
      {\fontsize{9.5}{12.5}\selectfont #3\par}\vspace{2pt}
    \fi
    \sbox{\publinks}{#5}%
    \begin{minipage}[t]{\dimexpr\linewidth-\wd\publinks-8pt\relax}
      \RaggedRight\hyphenpenalty=10000
      {\fontsize{9.2}{12}\selectfont\textcolor{muted}{#4}\par}
    \end{minipage}\hfill\usebox{\publinks}%
    \if\relax\detokenize{#1}\relax\else
      \par\vspace{4pt}
      {\fontsize{9.5}{12.5}\selectfont\color{accent}#1\par}
    \fi
    \par\vspace{3pt}
    {\fontsize{9.5}{12.5}\selectfont #6\par}
  \end{minipage}\par\vspace{9pt}
}

\newcommand{\presentationlink}[2]{%
  \href{#1}{\textcolor{accent}{\fontsize{9}{12}\selectfont #2}}%
}

\newcommand{\presentationentry}[4]{%
  \alignedline{{\fontsize{9.5}{13}\selectfont\textbf{#1}\textcolor{muted}{（#2）}}}{%
    \presentationlink{#3}{Paper}\hspace{0.5em}\textcolor{rulegray}{|}\hspace{0.5em}%
    \presentationlink{#4}{Slides}%
  }\par\vspace{4pt}%
}

\newcommand{\teachingentry}[3]{%
  \alignedline{\textbf{#1}}{#2}\par\vspace{2pt}
  {\fontsize{9.5}{12.5}\selectfont\color{muted}助教；#3\par}\vspace{6pt}
}

\newcommand{\honorentry}[2]{%
  \alignedline{#1}{#2}\par\vspace{4pt}
}

\begin{document}

\begin{center}
  {\fontsize{30}{35}\selectfont \textbf{李严}}\\[7pt]
  北京邮电大学 \\[4pt]
  {\fontsize{9.5}{13}\selectfont\color{muted}北京 \quad | \quad \href{mailto:liyanly@bupt.edu.cn}{liyanly@bupt.edu.cn}
  \quad | \quad \href{https://scholar.google.com/citations?user=LZgGl1kAAAAJ\&hl=zh-CN}{Google Scholar}}
\end{center}

\section*{教育背景}
\educationentry
  {北京邮电大学}
  {管理科学与工程硕士}
  {导师：\href{https://teacher.bupt.edu.cn/liyawen}{李雅文}；联合导师：\href{https://teacher.bupt.edu.cn/yeguanhua}{叶冠华}}
  {2025年9月至今}
  {北京}

\vspace{10pt}
\educationentry
  {电子科技大学}
  {金融学、电子信息工程双学士学位}
  {毕业论文导师：\href{https://www.mgmt.uestc.edu.cn/mgmt/jsnry.jsp?urltype=news.NewsContentUrl\&wbtreeid=1020\&wbnewsid=12383}{刘波}}
  {2021年9月至2025年7月}
  {成都}

\section*{论文与手稿}
\pubtopic{多模态医学影像分析}
  {我将临床语言和解剖先验融入影像模型，研究如何检索符合临床查询的图像，并定位、分割相关病灶。}

\pubentry
  {MedFuse: A Multi-source Data Fusion Framework for Diabetic Retinopathy Lesion Segmentation}
  {Jiantong Du, \textbf{Yan Li}, Yawen Li, Liwen Liao, Zhihao Zhao, and \corrauth{Guanhua Ye}.}
  {\textit{Frontiers of Computer Science} (\textbf{FCS}), 2026.}
  {\presentationlink{https://github.com/leeyan0612/CV/blob/main/papers/medfuse-fcs-2026.pdf}{Paper}}
  {与北京大学第一医院眼科医生合作，主导基于血管先验的双编码器融合。在 DDR 和 IDRiD 上，5 种骨干网络的 mIoU 和 mAUPR 均得到提升。}

\pubentry
  {Learning to ground compositional medical queries in retinal images}
  {Yawen Li, \textbf{Yan Li}, Jingbin Niu, Yingxia Shao, and \corrauth{Guanhua Ye}.}
  {\textit{Information Fusion}（投稿中）。}
  {\presentationlink{https://github.com/leeyan0612/CV/blob/main/papers/refine-retinal-grounding.pdf}{Paper}\hspace{0.5em}\textcolor{rulegray}{|}\hspace{0.5em}\presentationlink{https://github.com/leeyan0612/REFINE}{Code}}
  {主导组合查询建模与动态加权训练，实现眼底图像检索与病灶证据定位。在 3 个数据集和 7 类视觉语言模型上验证。}

\pubtopic{医疗 AI 评测与可靠性}
  {我研究医疗与多模态 AI 的评测可靠性，关注子群差异和评分规则如何掩盖模型失效。}

\pubentry[%
  \textbf{Outstanding Paper Award}（Responsible Multimedia 主题）。\par
  \textbf{Oral}（2026年11月，里约热内卢）：入选主会 \textbf{168} 篇 Oral 论文之一（共 \textbf{7,103} 篇送审论文）。]
  {Beyond Metadata: CAPRA for Hidden Subgroup Analysis under Missing Metadata in Medical Imaging}
  {Yawen Li, \textbf{Yan Li}, Zhe Xue, Yingxia Shao, Meiyu Liang, and \corrauth{Guanhua Ye}.}
  {\textit{ACM International Conference on Multimedia} (\textbf{ACM MM}), 2026.}
  {\presentationlink{https://github.com/leeyan0612/CV/blob/main/papers/capra-acm-mm-2026.pdf}{Paper}\hspace{0.5em}\textcolor{rulegray}{|}\hspace{0.5em}\presentationlink{https://github.com/leeyan0612/CV/blob/main/posters/capra-acm-mm-2026.pdf}{Poster}\hspace{0.5em}\textcolor{rulegray}{|}\hspace{0.5em}\presentationlink{https://github.com/leeyan0612/CV/blob/main/slides/capra-acm-mm-2026.pdf}{Slides}}
  {主导缺失元数据条件下的影像子群推断与概率校准。将推断结果用于鲁棒训练后，15 组学习方法与数据集组合中有 14 组提升了最差子群准确率。}

\pubentry
  {What Does Your Short-Answer VQA Score Actually Measure? Evaluator-Dependent Instability in Multimodal Short-Answer Benchmarks}
  {Guanhua Ye, Jingbin Niu, \textbf{Yan Li}, Meiyu Liang, Zhe Xue, Yingxia Shao, and \corrauth{Yawen Li}.}
  {\textit{ACM International Conference on Multimedia} (\textbf{ACM MM}), 2026.}
  {\presentationlink{https://github.com/leeyan0612/CV/blob/main/papers/vqa-evaluator-acm-mm-2026.pdf}{Paper}}
  {参与研究问题设计与论文写作。研究对逾 3.7 万条被判错的答案进行语义复核，发现 ST-VQA 上被判错答案中语义可接受的比例最高达 47.7\%。}

\clearpage
\section*{论文与手稿（续）}
\pubtopic{表征学习与预测}
  {我研究如何保留结构化数据的语义，并平衡异构特征，以改善下游预测。}

\pubentry
  {Rethink Representation Learning for Questionnaire Data}
  {Guanhua Ye, Jifeng He, \textbf{Yan Li}, Junping Du, Zhe Xue, Yingxia Shao, Meiyu Liang, and \corrauth{Yawen Li}.}
  {\textit{AAAI Conference on Artificial Intelligence} (\textbf{AAAI}), 2026.}
  {\presentationlink{https://github.com/leeyan0612/CV/blob/main/papers/questionnaire-representation-aaai-2026.pdf}{Paper}}
  {参与论文写作与修订。SemantiQ 利用检索增强生成保留问卷题目与回答的语义，在 3 个问卷数据集上的准确率和宏 F1 均优于基线。}

\pubentry[\textbf{Outstanding Paper Award}，第三届数据智能与管理学术会议（2023年）。]
  {FeBT: A Feature Balancing Transformer for Corporate ESG Forecasting}
  {Yawen Li, Mengyu Zhuang, \corrauth{Guanhua Ye}, \textbf{Yan Li}, Junheng Wang, Jinyi Zhou, and Pengfei Zhang.}
  {\textit{IEEE Transactions on Knowledge and Data Engineering} (\textbf{TKDE}), 2025.}
  {\presentationlink{https://github.com/leeyan0612/CV/blob/main/papers/febt-tkde-2025.pdf}{Paper}}
  {参与研究问题设计、第一阶段模型搭建及论文写作与修订。FeBT 使用掩码自编码器平衡输入 Transformer 的异构特征，并在真实企业数据上改善了 ESG 评分和评级预测。}

\section*{学术会议报告}
\presentationentry{Auditing Historical Clinical Queries with Missing Availability Time}{CCF BigData，2026}
  {https://github.com/leeyan0612/CV/blob/main/papers/chronoguard-ehr-ccf-bigdata-2026.pdf}
  {https://github.com/leeyan0612/CV/blob/main/slides/auditing-historical-clinical-queries-ccf-bigdata-2026.pdf}
\presentationentry{Slot-Based Evidence Routing for Budgeted Multi-Agent Coordination}{CINT，2026}
  {https://github.com/leeyan0612/CV/blob/main/papers/slot-based-evidence-routing-cint-2026.pdf}
  {https://github.com/leeyan0612/CV/blob/main/slides/slot-based-evidence-routing-cint-2026.pdf}
\presentationentry{Trust-free Personalized Decentralized Learning}{IEEE BigComp，2026}
  {https://github.com/leeyan0612/CV/blob/main/papers/trust-free-decentralized-learning-bigcomp-2026.pdf}
  {https://github.com/leeyan0612/CV/blob/main/slides/trust-free-personalized-decentralized-learning-bigcomp-2026.pdf}

\section*{教学经历}
\institutionheader{北京邮电大学}
\teachingentry{大数据与智能管理}{2025年秋季、2026年秋季}{共同设计面向低年级硕士和博士生的课程，班级规模为 130 人。}

\section*{荣誉与奖励}
\institutionheader{北京邮电大学}
\honorentry{一等奖学金}{2025、2026年}

\vspace{4pt}
\institutionheader{电子科技大学}
\honorentry{荣誉研究证书}{2025年}
\honorentry{经济与管理学院“学术之星”（唯一获奖者）}{2024年}
\honorentry{优秀学生奖学金}{2024年}
\honorentry{模范学生奖学金}{2023年}

\end{document}
